% fixture: amsart/amsart-maketitle
% tier: T3
% engine: document
% covers: amsart's \maketitle: the title in capitals with \@upprep's spaces, two authors \MakeUppercase'd, the abstract (\abstractbox, \@setabstract), \date, \subjclass and \keywords as footnotes without marks, \topskip 42pt on the first page, \ps@firstpage; the addresses (\interlinepenalty\@M) and emails at the end of the document
% issue: #360
\documentclass{amsart}
\begin{document}
\title{On the Distribution of Lattice Points in Large Convex Regions of the Plane}
\author{Ann Author}
\address{Department of Mathematics, Some University \\ Some City}
\email{ann@example.org}
\author{Bob Writer}
\address{Institute for Advanced Things, a Very Long Name of an Institute That Goes On and On Past the End of a Line}
\email{bob@example.org}
\date{May 2026}
\subjclass[2010]{Primary 11P21}
\keywords{lattice points, convex regions}

\begin{abstract}
We count the lattice points in large convex regions of the plane whose boundary has positive curvature, and give an error term that improves on the classical one by a power of the logarithm.

A second paragraph of the abstract, indented by the normal paragraph indentation of the class.
\end{abstract}

\maketitle

\section{Introduction}
Let $D$ be a convex region of the plane. The number of points of $\mathbb{Z}^2$ in the dilate $tD$ is the area of $tD$ up to an error term, and the problem is to bound that term as $t$ grows.

\section{Results}
Our main result is an improvement of the exponent for regions with smooth boundary.
\end{document}
